% !TeX root = ../../Book3.tex
\section{圆盘上的 Poincaré 度量}
\label{b3:sec:1301}

Schwarz--Pick 的分母给出圆盘内的局部长度单位。把这个单位沿曲线积分，再在曲线中取下确界，就得到距离；分式三角不等式则变成测地线上的普通加法。

\subsection{度量元}
\label{b3:subsec:130101}

圆盘的\term[poincare-duliang]{Poincaré 度量}[Poincaré metric]长度元为
\[
 \mathrm{d}s=\lambda_{\mathbb D}(z)|\mathrm{d}z|,
 \qquad \lambda_{\mathbb D}(z)=\frac2{1-|z|^2}.
\]
分段 \(C^1\) 路径的长度为 \(L(\gamma)=\int\lambda_{\mathbb D}(\gamma(t))|\gamma'(t)|\,\mathrm{d}t\)。本卷固定曲率为 \(-1\) 的归一化，不能删掉因子2而仍沿用下面的曲率和距离式。

\ProseParagraphBreak
密度在内部光滑且正，趋近边界时发散。内部小紧盘上它与欧氏密度有统一双边界，所以给相同局部拓扑；发散则把有限欧氏距离的边界推到无限内在距离。

\ProseParagraphBreak
度量元不是点之间的距离。它先给分段光滑路径长度 \(\ell(\gamma)=\int\lambda(\gamma)|\gamma'|\,\mathrm{d}t\)，再对连接两点的全部路径取下确界。因为 \(\lambda\) 在每个紧圆盘上有正上下界，局部长度与 Euclid 长度可比较；因此不同点的距离为正，且度量拓扑与原来拓扑相同。对圆盘中远近的比较不能仅把两个端点的密度平均，因为最短路径可能绕开密度大的区域。

\ProseParagraphBreak
本章归一化曲率为 \(-1\)，故密度是 \(2/(1-|z|^2)\)。也有文献使用 \(1/(1-|z|^2)\)，所有长度减半，曲率成为 \(-4\)。换常数并不改变测地线形状，但会改变距离、面积和闭曲线长度中的系数。后面圆环核心长度与模量的乘积为 \(\pi\)，也是这一归一化的结果。

\subsection{Möbius 不变性}
\label{b3:subsec:130102}

圆盘自同构满足 \(\lambda_{\mathbb D}(Tz)|T'(z)|=\lambda_{\mathbb D}(z)\)，已由第6章的导数恒等式得到。换元使每条路径长度不变，因此也保持路径距离。这个不变性把任意点的计算归到0。

\ProseParagraphBreak
对一般全纯圆盘映射只有不等式 \(\lambda_{\mathbb D}(f(z))|f'(z)|\le\lambda_{\mathbb D}(z)\)。于是每条像路径长度不超过原路径长度。若在一点等号且映射非常值，Schwarz--Pick 刚性使它成为自同构；长度不变的全局对称由局部导数等号识别。

\subsection{Poincaré 距离}
\label{b3:subsec:130103}

\begin{theorem}\label{b3:thm:13-disk-distance}
Poincaré 路径距离为
\[
 \mathrm d_{\mathbb D}(z,w)=2\operatorname{arctanh}\delta(z,w)
 =\log\frac{1+\delta(z,w)}{1-\delta(z,w)}.
\]
\end{theorem}
\begin{proof}
先令一端为0，另一端模为 \(r\)。对任意连接路径，记径向函数 \(u(t)=|\gamma(t)|\)，几乎处处 \(|u'|\le|\gamma'|\)。故其长度至少为 \(\int 2|u'|/(1-u^2)\,\mathrm{d}t\)，又至少为原函数 \(2\operatorname{arctanh}u\) 的端点差 \(2\operatorname{arctanh}r\)。径向线段达到此值。再用自同构把 \(w\) 移到0，像点模为 \(\delta(z,w)\)，不变性给公式。
\end{proof}
密度正使不同点距离正；路径拼接给三角不等式，反向给对称。第6章分式三角式经反双曲正切后也直接验证相同结论。

\begin{proposition}[距离的双曲余弦形式]\label{b3:prop:13-cosh-distance}
对 \(z,w\in\mathbb D\)，有
\[
 \cosh d(z,w)=1+\frac{2|z-w|^2}{(1-|z|^2)(1-|w|^2)}.
\]
若两条从0出发的径向线夹角为 \(\theta\)，端点到0的距离为 \(a,b\)，端点间距离为 \(c\)，则
\[
 \cosh c=\cosh a\cosh b-\sinh a\sinh b\cos\theta.
\]
\end{proposition}
\begin{proof}
令 \(q=\delta(z,w)\)，\(d=2\operatorname{artanh}q\)，故 \(\cosh d=(1+q^2)/(1-q^2)\)。再用
\[
 |1-\bar wz|^2-|z-w|^2=(1-|z|^2)(1-|w|^2)
\]
化简即得第一式。第二式取 \(z=r\)、\(w=s e^{\mathrm{i}\theta}\)，则分子为 \(r^2+s^2-2rs\cos\theta\)。由 \(r=\tanh(a/2)\)、\(s=\tanh(b/2)\)，\(\cosh a=(1+r^2)/(1-r^2)\)、\(\sinh a=2r/(1-r^2)\)，代入后逐项整理得所写公式。
\end{proof}
因共形度量保持夹角，\(\theta\) 也是该三角形在0的内角。把任一顶点用圆盘自同构送到0，便得到一般双曲三角形在该顶点的余弦关系。这个推导只用已经建立的距离公式，无需先接受一般曲面的三角公式。

\subsection{Schwarz–Pick 压缩性}
\label{b3:subsec:130104}

\begin{proposition}\label{b3:prop:13-disk-contraction}
圆盘全纯自映射压缩 Poincaré 距离。不同两点之间距离取等，或一个点的密度导数取等，都迫使映射是自同构。
\end{proposition}
\begin{proof}
沿任意源路径积分 Schwarz--Pick 导数不等式，得到像长度不大于源长度。目标距离不大于像长度，再对源路径取下确界。等号的两点形式也可直接对严格递增的 \(2\operatorname{arctanh}\) 使用伪双曲等号刚性；导数形式在第6章已证明。
\end{proof}
距离压缩比较的是域的内在单位。\(z\mapsto(z+a)/(1+az)\) 可有欧氏导数大于1，但它处处保持这里的长度。

\begin{proposition}[固定像值的紧致控制]\label{b3:prop:13-anchored-maps}
设 \(X\) 为圆盘型曲面，\(K\subset X\) 紧。全纯映射族 \(f:\mathbb D\to X\)、\(f(0)\in K\)，在每个闭子盘的像都包含于 \(X\) 的一个紧集，并且每个序列有局部一致收敛到全纯映射的子列。
\end{proposition}
\begin{proof}
第四节的定理\ref{b3:thm:13-hyperbolic-schwarz-pick}给 \(d_X(f(z),f(0))\le d_{\mathbb D}(z,0)\)。固定 \(|z|\le r<1\)，右边有统一上界。用有限个球覆盖 \(K\)，并扩大半径，像落在有限个紧闭球的并中；紧闭球性质由本章后面独立的覆盖距离证明给出。同样的压缩使族在每个子盘等度连续。在每个子盘选可数稠密点，利用紧目标逐点抽取并作对角子列。给定误差，等度连续性提供统一小邻域，再用有限网格覆盖子盘；有限网格点的收敛和三角不等式使子列在整个子盘一致 Cauchy。对可数圆盘穷竭再作对角抽取，得到局部一致极限 \(F\)。在任一点附近，把 \(F\) 的像限制到目标小坐标盘；一致收敛使子列充分晚也落在该盘中，其坐标表达的全纯极限由 Weierstrass 定理全纯。因此 \(F\) 是全纯曲面映射。
\end{proof}
固定0的像落在紧集的条件阻止整族向曲面端部逃走。在三穿孔球面上，它给省略值函数的一种紧致控制，但若基点像趋向穿孔，就可能只有球面值的退化极限。第十四章采用球面正规定义，能同时容纳这类退化；本命题则给非退化的曲面值极限。这两种紧致性陈述的目标不同，使用前应当检查需要哪一种极限。
